\documentclass[a4,11pt]{aleph-notas}
% Se puede ver la documentación aquí:
% https://github.com/alephsub0/LaTeX_aleph-notas

% -- Paquetes adicionales
\usepackage{enumitem}
\usepackage{url}
\usepackage{array}
\usepackage{booktabs}
\usepackage{longtable}
\usepackage{aleph-comandos}
\hypersetup{
    urlcolor=blue,
    linkcolor=blue,
}


% -- Datos
\institucion{Facultad de Ciencias Exactas, Naturales y Ambientales}
\carrera{Catálogo STEM}
\asignatura{Matemática III}
\tema{Plan tema 14: Cambio de variable en integrales triples}
\autor{Andrés Merino}
\fecha{Periodo 2026-1}

\logouno[0.16\textwidth]{Logos/logoPUCE_04_ac}
\definecolor{colortext}{HTML}{1957d1}
\definecolor{colordef}{HTML}{1957d1}
\fuente{montserrat}


\begin{document}

\encabezado

%%%%%%%%%%%%%%%%%%%%%%%%%%%%%%%%%%%%%%%%
\section*{Resultado de Aprendizaje}
%%%%%%%%%%%%%%%%%%%%%%%%%%%%%%%%%%%%%%%%

%%%%%%%%%%%%%%%%%%%%%%%%%%%%%%%%%%%%%%%%
\subsection*{RdA de la asignatura:}
\begin{itemize}[leftmargin=*]
    \item \textbf{RdA 1:} Identificar conceptos, teoremas y propiedades de funciones vectoriales, derivadas parciales e integral vectorial.
    \item \textbf{RdA 2:} Calcular derivadas parciales e integrales con asistentes computacionales.
    \item \textbf{RdA 3:} Aplicar Cálculo Vectorial para resolver problemas reales.
\end{itemize}

%%%%%%%%%%%%%%%%%%%%%%%%%%%%%%%%%%%%%%%%
\subsection*{RdA del tema:}
\begin{itemize}[leftmargin=*]
    \item Aplicar la fórmula de cambio de variable a integrales triples mediante el jacobiano de la transformación.
    \item Utilizar coordenadas cilíndricas para calcular integrales triples sobre sólidos con simetría cilíndrica.
    \item Utilizar coordenadas esféricas para calcular integrales triples sobre sólidos con simetría esférica.
    \item Seleccionar el sistema de coordenadas más adecuado según la geometría de la región de integración.
\end{itemize}

%%%%%%%%%%%%%%%%%%%%%%%%%%%%%%%%%%%%%%%%
\section*{Introducción}
%%%%%%%%%%%%%%%%%%%%%%%%%%%%%%%%%%%%%%%%

\paragraph{Pregunta inicial:}
Para calcular el volumen de una pelota de básquetbol (una esfera), un ingeniero plantea una integral triple en coordenadas rectangulares y obtiene límites de integración muy complicados. ¿Existirá una manera de describir la esfera con límites más simples que hagan el cálculo directo?

%%%%%%%%%%%%%%%%%%%%%%%%%%%%%%%%%%%%%%%%
\section*{Desarrollo}
%%%%%%%%%%%%%%%%%%%%%%%%%%%%%%%%%%%%%%%%

%%%%%%%%%%%%%%%%%%%%%%%%%%%%%%%%%%%%%%%%
\subsection*{Actividad 1: Coordenadas cilíndricas}
Se introduce el cambio de variable a coordenadas cilíndricas como extensión natural de las polares al espacio, y se aplica a sólidos con simetría de revolución alrededor del eje $z$.

\paragraph{¿Cómo lo haremos?}
\begin{itemize}[leftmargin=*]
    \item \textbf{Motivación:} mostrar que un cilindro o un paraboloide de revolución tiene una descripción mucho más simple en coordenadas cilíndricas que en rectangulares; conectar con las coordenadas polares del tema anterior.
    \item \textbf{Clase magistral:} se presenta la transformación $(r,\theta,z)\mapsto(r\cos\theta,r\sin\theta,z)$; se calcula el jacobiano $r$; se identifican los tipos de regiones convenientes (cilindros, conos, paraboloides); se aplica la fórmula de cambio de variable. Se usa el resumen \href{https://andres-merino.github.io/MatematicaIII-05-N0051/01Resumen/Resumen14.pdf}{Resumen14.pdf}.
    \item \textbf{Implementación computacional:} visualizar sólidos en coordenadas cilíndricas y calcular integrales triples con un asistente computacional.
    \item \textbf{Resolución de ejercicios:} calcular integrales triples en coordenadas cilíndricas usando \href{https://andres-merino.github.io/MatematicaIII-05-N0051/04Ejercicios/Ejercicios14.pdf}{Ejercicios14.pdf}.
\end{itemize}

\paragraph{Verificación de aprendizaje:}
\begin{itemize}[leftmargin=*]
    \item Calcule el volumen del sólido acotado por $z=r^2$ y $z=4$ usando coordenadas cilíndricas.
    \item ¿Por qué aparece el factor $r$ en la integral en coordenadas cilíndricas?
    \item ¿Qué regiones tridimensionales se describen con mayor comodidad en coordenadas cilíndricas? Dé dos ejemplos.
\end{itemize}

%%%%%%%%%%%%%%%%%%%%%%%%%%%%%%%%%%%%%%%%
\subsection*{Actividad 2: Coordenadas esféricas}
Se introduce el cambio de variable a coordenadas esféricas y se aplica al cálculo de integrales sobre sólidos con simetría esférica.

\paragraph{¿Cómo lo haremos?}
\begin{itemize}[leftmargin=*]
    \item \textbf{Motivación:} retomar la pregunta inicial; mostrar que una esfera de radio $a$ se describe simplemente como $0\leq\rho\leq a$, $0\leq\theta\leq 2\pi$, $0\leq\phi\leq\pi$ en coordenadas esféricas.
    \item \textbf{Clase magistral:} se presenta la transformación $(\rho,\theta,\phi)\mapsto(\rho\sin\phi\cos\theta, \rho\sin\phi\sin\theta, \rho\cos\phi)$; se calcula el jacobiano $\rho^2\sin\phi$; se identifican los tipos de regiones convenientes (esferas, casquetes, conos con vértice en el origen); se aplica la fórmula. Se usa el resumen \href{https://andres-merino.github.io/MatematicaIII-05-N0051/01Resumen/Resumen14.pdf}{Resumen14.pdf}.
    \item \textbf{Resolución de ejercicios:} calcular integrales triples en coordenadas esféricas usando \href{https://andres-merino.github.io/MatematicaIII-05-N0051/04Ejercicios/Ejercicios14.pdf}{Ejercicios14.pdf}.
\end{itemize}

\paragraph{Verificación de aprendizaje:}
\begin{itemize}[leftmargin=*]
    \item Calcule el volumen de la esfera de radio $a$ usando coordenadas esféricas.
    \item ¿Por qué aparece el factor $\rho^2\sin\phi$ en la integral en coordenadas esféricas? Interprete cada factor.
    \item Describa en coordenadas esféricas el casquete esférico $\{(x,y,z): x^2+y^2+z^2\leq 4,\, z\geq 1\}$.
\end{itemize}

%%%%%%%%%%%%%%%%%%%%%%%%%%%%%%%%%%%%%%%%
\section*{Cierre}
%%%%%%%%%%%%%%%%%%%%%%%%%%%%%%%%%%%%%%%%

\paragraph{Tarea:}
Resolver del libro \href{https://catalogobiblioteca.puce.edu.ec/cgi-bin/koha/opac-detail.pl?biblionumber=83405&query_desc=su%3A%22An%C3%A1lisis%20vectorial%22}{Cálculo Vectorial de Colley}: sección 5.5, ejercicios 27--35 (impares).

\paragraph{Pregunta de investigación:}
\begin{enumerate}[leftmargin=*]
    \item ¿En qué aplicaciones físicas o de ingeniería se usan las coordenadas esféricas? Investigar al menos un ejemplo concreto.
    \item ¿Cómo se relacionan las coordenadas cilíndricas y esféricas? ¿Es posible pasar de unas a otras directamente sin pasar por las rectangulares?
\end{enumerate}

\paragraph{Para el próximo tema:}
Ver los videos:
\begin{itemize}[leftmargin=*]
    \item \href{https://youtu.be/_60sKaoRmhU?si=ZxPAXnud8_u64IQK}{Introduction to the line integral}.
\end{itemize}


\end{document}
